\chapter{Conclusion}
\label{chap:conclusion}

This report started from a practical observation: the efficiency of modern algorithms cannot be fully explained by classical textbook techniques alone. In particular, high-performance regular expression matching is no longer explained only by classical automata theory. Engines such as Hyperscan combine automata constructions, graph decompositions, SIMD string matchers, hashing, and verification heuristics. This makes them interesting from two complementary points of view. On the systems side, their performance depends on hardware-aware design choices. On the mathematical side, several of those choices lead to natural statistics of random regular expressions and to recursive distributional equations.

We first reviewed the concepts needed to make this connection precise: regular expressions and finite automata, the Glushkov construction, BST-like distributions on syntax trees, SIMD parallelism, analytic combinatorics, and Monte Carlo simulation. We then traced the main compilation and execution paths of Hyperscan. In the compilation phase, regular expressions are converted to Glushkov NFAs, decomposed by Violet, and organized by Rose into a graph of literals and sub-FAs. In the execution phase, literal matchers act as filters that trigger the remaining automata. This viewpoint suggests that the extracted literals are not merely implementation details: their lengths and interactions directly influence how much work is delegated to SIMD matchers and how much work remains for verification.

The FDR string matcher provided a concrete example of this phenomenon. Its simplified model can be understood as a multi-pattern extension of Shift-Or, where patterns are grouped into buckets, represented through masks, and scanned with parameters such as super-character size and stride. The experiments around this model indicate that an analysis based only on SIMD operations is incomplete. False positives, hashing, and exact verification also have to be accounted for. This gives a first family of future problems: build a faithful probabilistic model of FDR that includes bucket assignment, super-characters, stride, the hash table, and the induced verification cost. Such a model should be compared against measurements on real pattern sets and, when possible, against the current Vectorscan implementation.

The main mathematical part of the report focused on the length of the longest literal extracted from a random regular expression. Under the BST-like model, concatenation leads to sums, alternation leads to maxima, and Kleene star resets the contribution to zero. This gives the stochastic recurrence \eqref{eq:longest-literal-recurrence}. For its expectation $(y_n)$, we introduced deterministic lower and upper bounds, then studied in detail the tightened lower-bound recurrence $(z_n)$. The results show that $(z_n)$ grows as a power of $n$: first in the weaker form $z_n=n^{\alpha+o(1)}$, then as $z_n=\Theta(n^\alpha)$, and finally, under the stronger condition $u>1/2$, as $z_n\sim Cn^\alpha$, where $\alpha$ is determined by a transcendental equation involving the parameters $u$ and $v$.

The proof of the exact asymptotics is interesting because the maximum operator cannot be removed immediately. The initial values create possible early fluctuations, so the proof first establishes eventual monotonicity through a negative-variation argument. This leaves a useful open question even for the deterministic recurrence: determine the first index from which monotonicity holds, and remove or weaken the current parameter restrictions. A long-term goal is to document a repository of techniques applicable to recurrences with sums and maxima, as they appear in many other contexts.

The original stochastic recurrence appears to have a richer asymptotic behavior than the deterministic lower bound. Simulations suggest a pure power growth, but with an exponent larger than the one obtained from the tightened lower bound. The fixed-point approach developed in the last part of the report gives a possible route forward. After normalizing $Y_n$ by $n^\alpha$, any limiting distribution should satisfy the distributional equation \eqref{eq:limiting-equation}. A direct contraction argument is obstructed by the degenerate fixed point $\delta_0$ and by the fact that $T_\alpha$ does not preserve the mean. The mean-preserving recurrence \eqref{eq:mean-preserving-recurrence} addresses this by evolving both the distribution and the exponent. We proved an existence result for a fixed point under a bounded second-moment condition, using Tychonoff's fixed-point theorem.

Several directions remain open. The most important one is to complete the asymptotic analysis of the stochastic recurrence: prove that $y_n=n^{\alpha^\star+o(1)}$, identify or characterize $\alpha^\star$, and prove convergence of the mean-preserving recurrence.

Overall, the internship shows that Hyperscan is a productive source of theoretical questions. Its optimizations are practical, but many of them expose clean mathematical objects: random syntax trees, sums and maxima in recurrences, dependency graphs, and filtering/verification trade-offs. A complete average-case analysis of such an engine is still out of reach, but the longest-literal problem gives a tractable first step and a concrete path toward more realistic models of modern regular expression matching.
